\documentclass[11pt]{article}
\usepackage[margin=1in]{geometry}
\usepackage[T1]{fontenc}
\usepackage{lmodern}
\usepackage{microtype}
\usepackage{amsmath,amssymb}
\usepackage{graphicx}
\usepackage{booktabs}
\usepackage{siunitx}
\usepackage{xcolor}
\usepackage[numbers,sort&compress]{natbib}
\usepackage[hidelinks]{hyperref}
\usepackage{caption}
\usepackage{longtable}
\captionsetup{font=small,labelfont=bf}
\graphicspath{{figures/}}
\sisetup{group-separator={,},group-minimum-digits=4}
\newcommand{\NTracts}{9,794}
\newcommand{\NHU}{17.1}
\newcommand{\NPop}{41.8}
\newcommand{\NBldg}{19.2}
\newcommand{\NPoi}{2.0}
\newcommand{\RefCorr}{0.975}
\newcommand{\RefMedRatio}{1.05}
\newcommand{\RefGapShare}{1.9\%}
\newcommand{\RefGapRural}{4.7\%}
\newcommand{\RefGapUrban}{1.3\%}
\newcommand{\RefGapTribal}{4.3\%}
\newcommand{\RefGapNonTribal}{1.6\%}
\newcommand{\CovBldg}{90.0\%}
\newcommand{\CovPoi}{90.0\%}
\newcommand{\FlagBldg}{510}
\newcommand{\FlagPoi}{529}
\newcommand{\RTwoBldg}{0.93}
\newcommand{\RTwoPoi}{0.82}
\newcommand{\CovSpatRegMinBldg}{88.9\%}
\newcommand{\CovSpatRegMaxBldg}{93.2\%}
\newcommand{\CovMondRegMinBldg}{88.9\%}
\newcommand{\CovMondRegMaxBldg}{90.2\%}
\newcommand{\CovLoroBldg}{84.6\%}
\newcommand{\CovLoroMinBldg}{78.1\%}
\newcommand{\PrecBldg}{0.09}
\newcommand{\RecBldg}{0.26}
\newcommand{\PrevBldg}{1.9\%}
\newcommand{\LiftBldg}{5.1}
\newcommand{\PRAUCBldg}{0.10}
\newcommand{\PRAUCOlsBldg}{0.06}
\newcommand{\PRAUCRatioBldg}{0.08}
\newcommand{\PrecPoi}{0.58}
\newcommand{\PrevPoi}{24.0\%}
\newcommand{\DelThirtyBldg}{66\%}
\newcommand{\DelFiftyBldg}{96\%}
\newcommand{\DelThirtyPoi}{18\%}
\newcommand{\DelFiftyPoi}{56\%}
\newcommand{\CountyRho}{0.85}
\newcommand{\CountyN}{324}
\newcommand{\NFdrFamily}{25}
\newcommand{\PoiSviHi}{3.8\%}
\newcommand{\PoiSviLo}{8.9\%}
\newcommand{\PoiSviAor}{0.41}
\newcommand{\HFiveWildRural}{4.8}
\newcommand{\HFiveWildRuralCI}{[2.4, 9.4]}
\newcommand{\HFiveDroughtRural}{2.9}
\newcommand{\HFiveDroughtRuralCI}{[1.2, 7.1]}
\newcommand{\HFiveHeatTribal}{0.08}
\newcommand{\HFiveHeatTribalCI}{[0.0, 0.4]}
\newcommand{\HEAll}{44,945}
\newcommand{\HFourDprBldg}{-0.010}
\newcommand{\HFourDprBldgCI}{[-0.035, +0.016]}
\newcommand{\HFourDaucPoi}{+0.004}
\newcommand{\HFourDaucPoiCI}{[+0.001, +0.007]}
\newcommand{\GapMsft}{2.8\%}
\newcommand{\GapOsm}{1.3\%}
\newcommand{\GainBldg}{1.7}
\newcommand{\OsmShareA}{31\%}
\newcommand{\OsmShareC}{51\%}
\newcommand{\GapShareA}{7.0\%}
\newcommand{\GapShareB}{1.8\%}
\newcommand{\GapRuralA}{15.7\%}
\newcommand{\GapRuralC}{4.2\%}
\newcommand{\GapSviC}{5.3\%}
\newcommand{\DidRuralPP}{9.9}
\newcommand{\DidRuralCI}{[-15.3, -5.5]}
\newcommand{\DidRuralP}{0.003}
\newcommand{\DidTribalPP}{6.7}
\newcommand{\DidTribalCI}{[-14.3, -0.6]}
\newcommand{\DidTribalP}{0.061}
\newcommand{\HThreeRural}{$+$0.060 log points, 95\% CI 0.013--0.106}
\newcommand{\HThreeRuralP}{0.040}
\newcommand{\EOneLost}{179,000}
\newcommand{\EOneTracts}{186}
\newcommand{\EOneBexar}{181}
\newcommand{\EOneResidents}{889,000}
\newcommand{\EOneSviShare}{63\%}
\newcommand{\EOneSviAll}{34\%}
\newcommand{\EOneFlagAfter}{82\%}
\newcommand{\EOneFlagBefore}{2.7\%}
\newcommand{\EOneFlagOther}{4.5\%}
\newcommand{\EOneRecall}{83\%}
\newcommand{\EOneRefGapAfter}{99\%}
\newcommand{\EThreeSummary}{rural +0.034 (endpoint 2026-08-19; 95\% CI 0.003--0.065); rural +0.040 (excluding Bexar County; 95\% CI 0.008--0.072); tribal +0.035 (endpoint 2026-08-19; 95\% CI 0.004--0.067); tribal +0.036 (excluding Bexar County; 95\% CI 0.002--0.070)}
\newcommand{\ETwoText}{For rural vs.\ urban tracts the adjusted odds ratio of a reference gap is 3.86 [2.40, 6.22], but 1.06 [0.72, 1.57] for the map-trained flag and 0.83 [0.48, 1.44] for the reference-anchored flag; for tribal vs.\ non-tribal tracts the values are 4.17 [1.89, 9.23], 1.26 [0.38, 4.20] and 0.92 [0.17, 5.14]. \emph{Absorption} is visible in the expectation itself: relative to the reference, observed completeness in rural tracts is lower by -0.031 [-0.053, -0.009] log points (region FE), and the map-trained expectation is already lower by -0.025 [-0.053, 0.004], so it reproduces 80\% of the rural shortfall (64\% of the tribal one). \emph{Resolution} limits even the reference-anchored model: its tract-level residual SD is 0.21 log points (map-trained: 0.19), against 0.22 for a 20\% gap, and it flags only 2.3\% of tracts (precision 0.16, recall 0.20 for reference gaps; map-trained 0.09 and 0.26).}
\newcommand{\OrigHOne}{4.9 vs.\ 2.3}
\newcommand{\OrigHOneP}{0.0004}
\newcommand{\RobustText}{Across 7 specifications (primary; no housing-type covariates; negative-binomial GLM; 80\% and 95\% ranges; within-region calibration; excluding tracts with fewer than 50 housing units), 53 of 56 flag contrasts keep the significance status of the primary analysis. The lower POI flag rate in high-SVI tracts holds in every specification (risk difference -0.040 to -0.067). The exceptions are: no housing-type covariates, building svi: RD +0.059 [+0.025, +0.086]; negative-binomial GLM, POI rural: RD -0.010 [-0.024, +0.002]; 95\% range, POI tribal: RD -0.012 [-0.028, -0.002]. Most notably, dropping units-in-structure covariates makes building flags significantly more common in high-SVI tracts (aOR 2.91), consistent with absorption: mobile-home and multi-unit housing shares, which correlate with SVI, explain away part of the high-SVI shortfall in the primary model. Details are in Appendix Table~\ref{tab:robust}.}
\newcommand{\RobustTextShort}{Across 7 specifications (primary; no housing-type covariates; negative-binomial GLM; 80\% and 95\% ranges; within-region calibration; excluding tracts with fewer than 50 housing units), 53 of 56 flag contrasts keep the significance status of the primary analysis. The lower POI flag rate in high-SVI tracts holds in every specification (risk difference -0.040 to -0.067). The exceptions are: no housing-type covariates, building svi: RD +0.059 [+0.025, +0.086]; negative-binomial GLM, POI rural: RD -0.010 [-0.024, +0.002]; 95\% range, POI tribal: RD -0.012 [-0.028, -0.002]. Most notably, dropping units-in-structure covariates makes building flags significantly more common in high-SVI tracts (aOR 2.91), consistent with absorption: mobile-home and multi-unit housing shares, which correlate with SVI, explain away part of the high-SVI shortfall in the primary model. Full results are in \texttt{outputs/tables/robustness.csv}.}


\title{Invisible Exposure: Reference-Free Conformal Auditing of Open-Map Completeness Reveals Catch-Up, a Release Regression, and Hidden Climate-Hazard Exposure}
\author{Chinthan\thanks{Corresponding author: chinthanm27@gmail.com} \and [Co-authors]}
\date{}

\begin{document}
\maketitle

\begin{abstract}
Open map data now underpin disaster planning, insurance, and climate-risk models, yet where these maps are incomplete, and for whom, is rarely measured, because a complete reference map seldom exists.
We propose a reference-free audit that predicts how many buildings and points of interest (POIs) a census tract \emph{should} contain from neutral structural facts (housing stock, population, land area, jobs), wraps every prediction in a cross-fitted conformalized-quantile-regression range, and flags tracts whose mapped count falls below it.
Vulnerability, rurality, tribal and hazard variables never enter the model; they are joined only in a pre-registered equity audit.
Applied to \NTracts{} tracts in five U.S. climate-hazard regions and three Overture Maps snapshots (2023--2026), the ranges achieve \CovBldg{} (buildings) and \CovPoi{} (POIs) held-out coverage of a 90\% target, but coverage degrades to \CovLoroBldg{} when whole regions are held out, and is restored by within-region calibration.
Against an independent reference (FEMA/ORNL USA Structures), flags are \LiftBldg{}$\times$ more precise than chance, and a deletion test detects 50\% feature loss in \DelFiftyBldg{} of tracts.
Building maps caught up between 2023 and 2026, faster in rural and tribal tracts (rural gap rate $-$\DidRuralPP{} percentage points relative to urban).
However, the \emph{October 2026} release silently dropped about \EOneLost{} building footprints, almost all Microsoft-derived, in \EOneTracts{} tracts of San Antonio, Texas (\EOneResidents{} residents; \EOneSviShare{} high-vulnerability tracts); the reference-free model flagged \EOneFlagAfter{} of them without any reference.
Conservatively counted missing structures in high wildfire-hazard tracts are \HFiveWildRural{}$\times$ more common per housing unit in rural than urban tracts.
Finally, we show a general limitation with practical consequences: rural and tribal building gaps that are 3--4$\times$ more prevalent in the reference are invisible to tract-level flags, because a map-trained expectation absorbs most of the group-level shortfall and because structural prediction noise exceeds the size of typical gaps; such models detect large idiosyncratic losses, not modest systematic ones.
Code and pinned data versions reproduce every figure.
\end{abstract}

\noindent\textbf{Keywords:} open map data; Overture Maps; data completeness; conformal prediction; environmental justice; climate hazards; algorithmic auditing

\section{Introduction}

Open geospatial data have become infrastructure.
OpenStreetMap, machine-learned building footprints, and their merger in Overture Maps \citep{overture} are used to count exposed buildings in wildfire and heat models, to route evacuations, and to locate services.
When a map omits buildings or businesses, the omission propagates: structures that are not mapped are not counted as exposed, and communities whose places are missing appear to need less.
Completeness of volunteered and open data is known to vary with affluence and urbanity \citep{haklay2010,mashhadi2013,thebaultspieker2018,herfort2023}, and road networks have been shown to approach completeness globally \citep{barringtonleigh2017}.
What remains hard is to estimate completeness \emph{where no complete reference exists}, with honest uncertainty, and to tie it to who bears climate risk.

We contribute a reference-free, uncertainty-aware audit and apply it to a pre-registered equity question.
The design follows one rule: the model that predicts how many features a tract should contain sees only neutral structural facts, never vulnerability, rurality, tribal status or hazard. These variables are joined afterwards, so the model cannot learn existing disparities as ``normal'' through them.
Our contributions are:
\begin{enumerate}
\item \textbf{A conformal completeness model.} Cross-fitted LightGBM quantile regression with conformalized quantile regression (CQR) \citep{romano2019} gives each tract a calibrated 90\% range for its expected building and POI count. We show that random and spatial-block cross-fitting attain nominal coverage, that leave-one-region-out transfer does not, and that Mondrian (within-region) calibration restores region-conditional coverage.
\item \textbf{External validation} against FEMA/ORNL USA Structures \citep{usastructures} and Census Business Patterns \citep{cbp}, with a controlled deletion test of detection power.
\item \textbf{A longitudinal audit} across three Overture snapshots spanning 2023--2026, measuring whether gaps in disadvantaged tracts closed, and documenting a large, silent release regression that the reference-free model detects.
\item \textbf{Hidden exposure}: missing structures in high heat, wildfire and drought hazard tracts, audited by social vulnerability (SVI) \citep{flanagan2011}, rurality and tribal status with a frozen pre-registration and false-discovery-rate control \citep{bh1995}.
\item \textbf{A detectability result}: tract-level flags from a self-referential expectation miss gaps shared by a whole group of similar tracts. Two mechanisms explain this. The expectation absorbs most of the group shortfall, which mirrors label-bias problems in algorithmic auditing \citep{obermeyer2019}. And structural prediction noise exceeds the size of typical gaps, which even a reference-trained expectation cannot overcome.
\end{enumerate}

\section{Related work}

\paragraph{Open-map completeness.}
Early OpenStreetMap quality studies compared volunteered data with authoritative maps and found that completeness tracked population density and affluence \citep{haklay2010}. Contributor behaviour explains part of this geography: biases are present from contributors' first edits \citep{thebaultspieker2018}, and POI coverage is lower in deprived areas \citep{mashhadi2013}. At global scale, \citet{herfort2023} estimated building completeness with regression on socio-economic covariates and found that large parts of the urban world remain under-mapped, with inequality between high- and low-income regions. Machine-learned footprints from imagery have since changed the picture in many countries, and Overture Maps fuses them with OpenStreetMap and other sources. To our knowledge, no study has audited Overture's completeness with calibrated uncertainty, across releases, or in relation to climate-hazard exposure.

\paragraph{Conformal prediction and spatial validation.}
Split and cross conformal methods give finite-sample marginal coverage under exchangeability \citep{vovk2005,lei2018,barber2021}; CQR adapts interval width to heteroscedastic noise \citep{romano2019}. Coverage can fail under covariate or distribution shift \citep{tibshirani2019,gibbs2021}. Spatial data violate exchangeability across space, and random cross-validation overstates spatial transfer \citep{roberts2017,ploton2020}. We bring these strands together for map auditing and test them empirically.

\paragraph{Equity audits and label bias.}
Audits of algorithms that learn from historical outcomes show that a target that encodes unequal access reproduces it \citep{obermeyer2019}. Our completeness model faces the same risk, because it learns from the very map it audits. We therefore quantify how much group-structured incompleteness the model absorbs.

\section{Data}

\paragraph{Study area.} We use the five study regions of the \emph{Bias Bounty: Mapping Equity} challenge \citep{biasbounty}: Maricopa County and surroundings (AZ), the Butte--Shasta--Tehama fire corridor (northern CA), eastern Oklahoma, south-central Texas, and eastern Washington. Together they contain \NTracts{} 2020 census tracts, \NHU{} million housing units and \NPop{} million residents. Each region is the dissolve of its tracts, so map extracts never cut a tract.

\paragraph{Map under audit.} Overture Maps buildings and places, release \texttt{2026-08-19.0} (primary; \NBldg{} million buildings and \NPoi{} million places), plus \texttt{2026-10-21.0} and an archived GeoParquet copy of \texttt{2023-07-26-alpha.0} for the longitudinal analysis. Each feature is reduced to its bounding-box centre and assigned to the containing tract. For buildings we count footprints of at least 450\,sq\,ft (\SI{41.8}{m^2}) to match the reference inclusion rule; footprint area is approximated from planar degrees scaled by $\cos(\text{latitude})$. We also record each feature's primary upstream source (OpenStreetMap, Microsoft ML Buildings, Google Open Buildings, Esri Community Maps; Meta, Microsoft and Foursquare for places).

\paragraph{Structural covariates (model inputs).} ACS 2020--2024 five-year counts \citep{acs}: housing units, population, households, group-quarters population, vacancy, units in structure (B25024), year built (B25034); land area; and LODES 2022 workplace jobs by NAICS sector \citep{lodes}. We use no worker demographics.

\paragraph{References (validation only).} USA Structures, ORNL's national inventory of structures over 450\,sq\,ft, carries a census block code, so tract counts need no spatial join. For POIs we allocate 2023 ZIP Business Patterns establishments to tracts by ZCTA land-area share, and use exact County Business Patterns totals at county level \citep{cbp}.

\paragraph{Audit strata (audit only).} CDC/ATSDR SVI 2022, the Climate Vulnerability Index, USDA RUCA 2020, tribal-area overlap, CDC EPHT summer extreme-heat days, USFS Wildfire Risk to Communities burn probability \citep{wrc}, and the U.S. Drought Monitor summer severity index, as distributed with the challenge \citep{biasbounty}.

\section{Methods}

\subsection{Expectation model with conformal ranges}
Let $y_i=\log(1+c_i)$ be the log count of mapped features in tract $i$ and $x_i$ its structural covariates. A guard in the code refuses any audit, reference, outcome or provenance column as an input. We fit LightGBM \citep{ke2017} quantile regressions for $\tau\in\{\alpha/2,\,0.5,\,1-\alpha/2\}$ with $\alpha=0.1$.
Predictions are \emph{cross-fitted}: tracts are partitioned into $K=10$ folds of whole \SI{50}{km} spatial blocks, and each tract is predicted by models that never saw its fold, so its own deficit cannot lower its expectation.
CQR conformity scores $s_i=\max\{\hat q_{\alpha/2}(x_i)-y_i,\; y_i-\hat q_{1-\alpha/2}(x_i)\}$ are computed out of fold. For fold $k$, the correction $\hat Q_k$ is the $\lceil(1-\alpha)(m+1)\rceil$-th smallest score among the $m$ tracts in the other folds. The range is $[\hat q_{\alpha/2}-\hat Q_k,\;\hat q_{1-\alpha/2}+\hat Q_k]$. In the Mondrian variant, scores are pooled within region.
A tract is \emph{flagged} when $y_i$ falls below the lower bound. Its conservative shortfall is $\max(\text{lower}-c_i,0)$ features, and its standardized deficit is $(\hat q_{0.5}-y_i)/(\hat q_{0.5}-\text{lower})$.

\subsection{Validation}
We compare random, spatial-block and leave-one-region-out (LORO) fold schemes, each with global or within-region calibration.
Against the references, a building gap is a tract whose Overture count is below 80\% of USA Structures. A POI gap is a tract whose Overture-to-establishment ratio is below half its region's median.
We report precision, recall, lift, ROC-AUC and PR-AUC, and compare with two baselines: a log-log OLS on housing units and land area, and the raw map-to-housing ratio.
In the deletion test we remove a share of a tract's mapped features from its observed count and record how often it is flagged. This is valid because the tract's observation never entered its own model or calibration fold.

\subsection{Pre-registered audit}
The project began with a pre-registration template (repository commit \texttt{807831b}, 2026-10-16) that stated five hypotheses and Mann--Whitney tests but left thresholds blank. Before any model output was cross-tabulated with any stratum, we completed and froze it (2026-11-09, commit \texttt{5d08b3e}), fixing thresholds and group definitions and specifying the regression-based contrasts below. Only reference-based Phase~1 descriptives had been seen, and these are disclosed in the pre-registration. We report the completed hypotheses as primary, and the original five hypotheses, tested exactly as worded with the blanks filled, in Appendix Table~\ref{tab:orig}.
\textbf{H1/H2}: flag rates differ between SVI Q4 and Q1, top and bottom CVI quartile, rural (RUCA $\geq 4$) and urban, and tribal and non-tribal tracts. We estimate a logistic regression with region fixed effects and county-clustered standard errors, and a risk difference with a county-cluster bootstrap (2,000 replicates).
\textbf{H3}: from 2023 to 2026-10, reference-relative completeness $\log[(c+1)/(r+1)]$ improves more in disadvantaged tracts. We use OLS with region fixed effects, plus a difference-in-differences of gap rates.
\textbf{H4}: provenance features improve gap detection (spatial-CV PR-AUC difference).
\textbf{H5}: \emph{hidden exposure} (conservative missing structures in top-tercile heat, wildfire or drought tracts, per 1,000 housing units) differs by group, as a rate ratio with cluster-bootstrap CI.
All H1, H2, H3 and H5 contrasts form one Benjamini--Hochberg family ($q=0.05$; \NFdrFamily{} tests). Analyses not in the pre-registration are labelled \emph{exploratory}.

\section{Results}

\subsection{How complete is the map?}
Overture's 2026-08 building layer agrees closely with USA Structures (log-scale Pearson $r=\RefCorr{}$; median ratio \RefMedRatio{}; Fig.~\ref{fig:scatter}). Only \RefGapShare{} of tracts fall below 80\% of the reference. Gaps are more common in rural (\RefGapRural{}) than urban (\RefGapUrban{}) tracts and in tribal (\RefGapTribal{}) than non-tribal (\RefGapNonTribal{}) tracts.
Under spatial cross-fitting the building model explains $R^2=\RTwoBldg{}$ of log-count variance, and the POI model $R^2=\RTwoPoi{}$ once job counts are included. The 90\% ranges cover \CovBldg{} and \CovPoi{} of held-out tracts, and \FlagBldg{} building and \FlagPoi{} POI tracts are flagged (Fig.~\ref{fig:maps}).

\begin{figure}[t]\centering
\includegraphics[width=\textwidth]{fig1_maps.png}
\caption{Standardized building deficit per tract (Overture 2026-08 against the conformal expectation) in the five study regions. Black outlines mark tracts below the 90\% conformal lower bound.}\label{fig:maps}
\end{figure}

\subsection{Calibration under spatial shift}
Random and spatial-block cross-fitting both reach nominal overall coverage, but per-region coverage under global calibration ranges from \CovSpatRegMinBldg{} to \CovSpatRegMaxBldg{} for buildings (Fig.~\ref{fig:cov}). Within-region calibration narrows the spread to \CovMondRegMinBldg{}--\CovMondRegMaxBldg{}.
When an entire region is held out, overall building coverage falls to \CovLoroBldg{}, with the worst region at \CovLoroMinBldg{}. Within-region calibration cannot help here, because the held-out region contributes no calibration scores.
In practice, a completeness model transferred to a new region needs some local calibration data to keep its guarantee.

\begin{figure}[t]\centering
\includegraphics[width=0.85\textwidth]{fig3_coverage_schemes.pdf}
\caption{Empirical coverage of 90\% conformal ranges by fold scheme and calibration. Dots: per region; diamonds: overall. Under leave-one-region-out folds the held-out region has no calibration scores, so the per-region option coincides with global calibration.}\label{fig:cov}
\end{figure}

\subsection{Do flags find real gaps?}
Against USA Structures, building flags have precision \PrecBldg{} and recall \RecBldg{} at a prevalence of \PrevBldg{}, a lift of \LiftBldg{}$\times$ (Table~\ref{tab:skill}). Their PR-AUC (\PRAUCBldg{}) exceeds the OLS (\PRAUCOlsBldg{}) and ratio (\PRAUCRatioBldg{}) baselines, while OLS has the higher ROC-AUC. The deletion test (Fig.~\ref{fig:del}) shows that removing 30\% of a tract's buildings triggers a flag \DelThirtyBldg{} of the time and removing 50\% \DelFiftyBldg{}; for POIs, whose expectation is noisier, the rates are \DelThirtyPoi{} and \DelFiftyPoi{}.
POI flags agree only moderately with the allocated establishment reference (precision \PrecPoi{}, prevalence \PrevPoi{}). Simple baselines do slightly better on this reference, because the land-area allocation of ZIP establishments is itself coarse. At county level, Overture places and Business Patterns establishments correlate at Spearman $\rho=\CountyRho{}$ ($n=\CountyN{}$).

\begin{table}[t]\centering\small
\caption{Flag skill against independent references (all tracts with a defined reference).}\label{tab:skill}
\begin{tabular}{llrrrrrrrr}
\toprule
Layer & Method & $n$ & Prev. & Flags & Prec. & Recall & Lift & ROC-AUC & PR-AUC \\
\midrule
Buildings & conformal model & 9,793 & 1.9\% & 510 & 0.094 & 0.264 & 5.1 & 0.685 & 0.099 \\
Buildings & OLS log-HU & 9,793 & 1.9\% & 510 & 0.075 & 0.209 & 4.0 & 0.727 & 0.060 \\
Buildings & map/housing ratio & 9,793 & 1.9\% & 510 & 0.090 & 0.253 & 4.9 & 0.463 & 0.081 \\
POIs & conformal model & 9,309 & 24.0\% & 495 & 0.584 & 0.130 & 2.4 & 0.671 & 0.415 \\
POIs & OLS log-HU & 9,309 & 24.0\% & 495 & 0.699 & 0.155 & 2.9 & 0.780 & 0.535 \\
POIs & map/housing ratio & 9,309 & 24.0\% & 495 & 0.675 & 0.150 & 2.8 & 0.764 & 0.514 \\
\bottomrule
\end{tabular}
\end{table}

\begin{figure}[t]\centering
\begin{minipage}{0.48\textwidth}\centering\includegraphics[width=\textwidth]{fig2_reference_scatter.png}
\caption{Overture vs.\ USA Structures buildings per tract; orange: flagged tracts.}\label{fig:scatter}\end{minipage}\hfill
\begin{minipage}{0.48\textwidth}\centering\includegraphics[width=\textwidth]{fig4_deletion.pdf}
\caption{Deletion test: share of tracts flagged after removing a share of their mapped features.}\label{fig:del}\end{minipage}
\end{figure}

\subsection{Pre-registered equity contrasts}
Table~\ref{tab:conf} lists all confirmatory contrasts.
\textbf{Buildings (H1, H2).} Model flag rates do not differ significantly by SVI, CVI, rurality or tribal status (all BH-adjusted $p>0.4$).
\textbf{POIs (H1, H2).} POI flags are \emph{less} common in high-SVI tracts (\PoiSviHi{} vs.\ \PoiSviLo{}; aOR \PoiSviAor{}), high-CVI tracts and rural tracts. The reference shows the same direction: Overture places per establishment are higher in high-SVI tracts. We find no evidence that Overture places under-cover vulnerable communities in these regions.
\textbf{Hidden exposure (H5).} The model attributes \HEAll{} conservatively missing structures to flagged tracts. Those in top-tercile wildfire-hazard tracts are \HFiveWildRural{}$\times$ more common per housing unit in rural than urban tracts (95\% CI \HFiveWildRuralCI{}), and in drought tracts \HFiveDroughtRural{}$\times$ (95\% CI \HFiveDroughtRuralCI{}). Hidden heat exposure is lower in tribal tracts (\HFiveHeatTribal{}$\times$).
\textbf{H4.} Provenance features did not improve detection of building gaps ($\Delta$PR-AUC \HFourDprBldg{}, 95\% CI \HFourDprBldgCI{}) and gave a marginal gain for POIs ($\Delta$ROC-AUC \HFourDaucPoi{}, 95\% CI \HFourDaucPoiCI{}). Tracts whose buildings come mainly from Microsoft ML footprints have twice the reference gap rate of OSM-dominated tracts (\GapMsft{} vs.\ \GapOsm{}). This is descriptive: once the model deficit is known, source mix adds little.

\begin{table}[t]\centering\small
\caption{Confirmatory contrasts (pre-registered). RD: risk difference (cluster-bootstrap 95\% CI). aOR: odds ratio adjusted for region (county-clustered 95\% CI). RR: hidden-exposure rate ratio. $\Delta$: difference in log completeness change (H3a) or in gap-rate change (H3b). $p_{BH}$: Benjamini--Hochberg adjusted.}\label{tab:conf}
\resizebox{\textwidth}{!}{\begin{tabular}{lllllrrc}
\toprule
Hyp. & Layer & Contrast & Group values & Estimate [95\% CI] & $p$ & $p_{BH}$ & Sig. \\
\midrule
H1 & Bldg & SVI Q4 vs Q1 & 6.0\% / 5.1\% & RD +0.009 [-0.010, +0.025]; aOR 1.19 & 0.278 & 0.535 &  \\
H1 & Bldg & CVI top vs bottom & 6.3\% / 4.5\% & RD +0.018 [-0.005, +0.035]; aOR 1.25 & 0.210 & 0.438 &  \\
H2 & Bldg & Rural vs urban & 5.0\% / 5.2\% & RD -0.002 [-0.022, +0.019]; aOR 1.06 & 0.765 & 0.822 &  \\
H2 & Bldg & Tribal vs non-tribal & 4.2\% / 5.3\% & RD -0.011 [-0.033, +0.013]; aOR 1.12 & 0.850 & 0.850 &  \\
H1 & POI & SVI Q4 vs Q1 & 3.8\% / 8.9\% & RD -0.051 [-0.071, -0.034]; aOR 0.41 & 0.000 & 0.000 & \checkmark \\
H1 & POI & CVI top vs bottom & 4.0\% / 7.5\% & RD -0.035 [-0.053, -0.019]; aOR 0.58 & 0.000 & 0.003 & \checkmark \\
H2 & POI & Rural vs urban & 3.7\% / 5.8\% & RD -0.021 [-0.034, -0.009]; aOR 0.56 & 0.001 & 0.003 & \checkmark \\
H2 & POI & Tribal vs non-tribal & 4.9\% / 5.4\% & RD -0.005 [-0.025, +0.012]; aOR 0.84 & 0.530 & 0.663 &  \\
H5 heat & Bldg & SVI Q4 vs Q1 & 1.41 / 1.03 & RR 1.37 [0.69, 3.54] & 0.313 & 0.559 &  \\
H5 heat & Bldg & Rural vs urban & 0.77 / 1.27 & RR 0.60 [0.18, 2.49] & 0.508 & 0.663 &  \\
H5 heat & Bldg & Tribal vs non-tribal & 0.11 / 1.29 & RR 0.08 [0.00, 0.36] & 0.001 & 0.003 & \checkmark \\
H5 wildfire & Bldg & SVI Q4 vs Q1 & 1.22 / 0.97 & RR 1.25 [0.72, 2.70] & 0.417 & 0.582 &  \\
H5 wildfire & Bldg & Rural vs urban & 3.19 / 0.67 & RR 4.79 [2.38, 9.41] & 0.001 & 0.003 & \checkmark \\
H5 wildfire & Bldg & Tribal vs non-tribal & 1.99 / 0.98 & RR 2.02 [0.64, 4.79] & 0.206 & 0.438 &  \\
H5 drought & Bldg & SVI Q4 vs Q1 & 0.86 / 1.04 & RR 0.83 [0.44, 3.65] & 0.776 & 0.822 &  \\
H5 drought & Bldg & Rural vs urban & 1.78 / 0.61 & RR 2.92 [1.23, 7.13] & 0.014 & 0.044 & \checkmark \\
H5 drought & Bldg & Tribal vs non-tribal & 1.26 / 0.75 & RR 1.67 [0.39, 4.67] & 0.419 & 0.582 &  \\
H3a & Bldg & SVI Q4 vs Q1 & 0.065 / 0.055 & $\Delta$ +0.010 [-0.054, +0.074] & 0.765 & 0.822 &  \\
H3a & Bldg & Rural vs urban & 0.119 / 0.059 & $\Delta$ +0.060 [+0.013, +0.106] & 0.011 & 0.040 & \checkmark \\
H3a & Bldg & Tribal vs non-tribal & 0.079 / 0.069 & $\Delta$ +0.036 [+0.002, +0.070] & 0.038 & 0.095 &  \\
H3a & Bldg & CVI top vs bottom & 0.086 / 0.050 & $\Delta$ +0.038 [-0.045, +0.121] & 0.373 & 0.582 &  \\
H3b & Bldg & SVI Q4 vs Q1 & -- & $\Delta$ -0.006 [-0.052, +0.059] & 0.789 & 0.822 &  \\
H3b & Bldg & CVI top vs bottom & -- & $\Delta$ -0.036 [-0.108, +0.047] & 0.340 & 0.567 &  \\
H3b & Bldg & Rural vs urban & -- & $\Delta$ -0.099 [-0.153, -0.055] & 0.001 & 0.003 & \checkmark \\
H3b & Bldg & Tribal vs non-tribal & -- & $\Delta$ -0.067 [-0.143, -0.006] & 0.022 & 0.061 &  \\
\bottomrule
\end{tabular}}
\end{table}

\begin{figure}[t]\centering
\includegraphics[width=0.75\textwidth]{fig7_hidden_exposure.pdf}
\caption{Hidden exposure rate ratios (H5) with cluster-bootstrap 95\% CIs; diamonds are significant after BH correction.}\label{fig:he}
\end{figure}

\subsection{Catch-up across releases (H3) and a silent regression}
Between 2023-07 and 2026-10 the study area gained \GainBldg{} million mapped buildings, and the OpenStreetMap share rose from \OsmShareA{} to \OsmShareC{}. The share of tracts below 80\% of the reference fell from \GapShareA{} (2023) to \GapShareB{} (2026-08) (Fig.~\ref{fig:catchup}).
Catch-up was faster where gaps had been largest: rural tracts went from \GapRuralA{} to \GapRuralC{}, a pre-registered difference-in-differences of $-$\DidRuralPP{} percentage points relative to urban (95\% CI \DidRuralCI{}; $p_{BH}=\DidRuralP{}$). Tribal tracts improved by $-$\DidTribalPP{} points relative to non-tribal (95\% CI \DidTribalCI{}; $p_{BH}=\DidTribalP{}$, not significant after FDR correction). Log-completeness gains were larger in rural tracts (\HThreeRural{}; $p_{BH}=\HThreeRuralP{}$). SVI and CVI contrasts were not significant.

The 2026-10 endpoint, however, hides a reversal.
Beginning with release \texttt{2026-09-23.0}, Overture dropped about \EOneLost{} footprints, almost all from Microsoft ML Buildings, from \EOneTracts{} tracts, \EOneBexar{} of them in Bexar County (San Antonio), TX. In the affected tracts the share below 80\% of the reference rose from 0\% to \EOneRefGapAfter{} (Fig.~\ref{fig:bexar}).
The affected tracts are home to \EOneResidents{} people, and \EOneSviShare{} of them are SVI-Q4 tracts, against \EOneSviAll{} study-wide. This regression is why the 2026-10 gap share among high-SVI tracts (\GapSviC{}) is higher than in 2026-08.
The reference-free model, re-run on the 2026-10 release with no reference input, flagged \EOneFlagAfter{} of the affected tracts, against \EOneFlagBefore{} on the 2026-08 release. Its recall for reference-confirmed gaps among them was \EOneRecall{}. Elsewhere the flag rate stayed at \EOneFlagOther{}.
In exploratory sensitivity analyses that end in 2026-08 or exclude Bexar County, the rural and tribal catch-up estimates remain positive with confidence intervals excluding zero: \EThreeSummary{} (Appendix Table~\ref{tab:e3}).

\begin{figure}[t]\centering
\includegraphics[width=\textwidth]{fig6_catchup.pdf}
\caption{Share of tracts below 80\% of USA Structures, by group and Overture snapshot. The 2026-10 uptick among high-SVI and high-CVI tracts is the Bexar County regression.}\label{fig:catchup}
\end{figure}

\begin{figure}[t]\centering
\includegraphics[width=0.6\textwidth]{fig8_bexar_regression.png}
\caption{Exploratory: buildings lost per tract between the 2026-08 and 2026-10 Overture releases around San Antonio, TX. Black outlines: tracts flagged by the reference-free model on the 2026-10 release.}\label{fig:bexar}
\end{figure}

\subsection{Normalization bias (exploratory)}
The reference shows building gaps 3--4$\times$ more often in rural and tribal tracts, yet the model's flags do not.
Fig.~\ref{fig:lenses} contrasts three outcomes: the map-trained flag; a flag from the same structural model trained on USA Structures counts (``reference-anchored''), raised when Overture falls below its lower bound; and the raw reference gap.
\ETwoText{}
Two mechanisms operate. First, \emph{absorption}: a model that learns from the map treats much of the group-wide under-mapping as the norm for structurally similar tracts. Second, \emph{resolution}: tract-level structural prediction noise is larger than a 20\% gap, so a per-tract lower bound cannot isolate modest shortfalls even when the expectation is trained on the reference. Per-tract flags therefore detect large, idiosyncratic losses such as the Bexar regression, but not modest, systematic ones. Null equity results from tract-level completeness flags should be read as uninformative about systematic disparity, not as evidence of its absence. Direct comparison with a reference, aggregated by group, is needed for that question.

\begin{figure}[t]\centering
\includegraphics[width=0.8\textwidth]{fig5_two_lenses.pdf}
\caption{Exploratory: adjusted odds ratios of a building gap by group under three lenses. Only the direct reference comparison reveals the rural and tribal disparity; tract-level flags from either expectation do not.}\label{fig:lenses}
\end{figure}

\subsection{Original pre-registered hypotheses}
Tested as originally worded (Appendix Table~\ref{tab:orig}): hidden exposure per 1,000 housing units is higher in high-SVI $\times$ high-heat than in low-SVI $\times$ low-heat tracts (mean \OrigHOne{}; $p_{BH}=\OrigHOneP{}$), although the medians of both groups are zero, because most tracts have no shortfall. Contrary to the original H2, rural or tribal high-drought tracts have \emph{lower} single-source dependence than other tracts. Contrary to the original H3, completeness improved slightly \emph{faster} in high-SVI tracts on the rank test; the regression contrast with region fixed effects (Table~\ref{tab:conf}) is not significant. Hidden exposure does not differ between high- and low-wildfire-hazard quartiles (original H5).

\subsection{Robustness}
\RobustText{}

\section{Discussion}
\paragraph{What the audit says about Overture.}
In these five regions, Overture's building layer is close to complete. The 2023--2026 period closed most rural and tribal building gaps, consistent with the large-scale import of machine-learned footprints and growth in OpenStreetMap. Places are not under-represented in vulnerable tracts relative to establishments. The largest single failure we observed was not a long-standing equity gap but a release regression, concentrated in high-vulnerability neighbourhoods of a major city, that a pinned-release user would not notice.

\paragraph{What the method contributes.}
Conformal ranges turn a completeness model into a monitor with an explicit false-alarm rate: about 5\% of tracts are flagged when nothing is wrong, and a 50\% loss is detected almost always. Because the model needs no reference, it can screen every new release; the Bexar case shows it working on a real event. The spatial-shift results argue for calibrating within each region, and caution against transferring a calibrated model to a new region without local data.

\paragraph{The detectability trap.}
Excluding vulnerability variables from the model keeps it from learning disparities \emph{through} those variables. It does not stop the model from learning them through structural correlates, such as low density and dispersed housing in rural and tribal areas. Nor does it give the model the resolution to see a 20\% shortfall in a single tract. We therefore recommend separating two uses. Reference-free conformal flags are a release-monitoring and triage tool with a known false-alarm rate. Equity conclusions about systematic gaps require a reference comparison aggregated by group, and should report the absorption share alongside any null result.

\paragraph{Limitations.}
USA Structures and the machine-learned footprints in Overture are both derived from imagery, so shared omissions (for example under tree canopy) would understate true gaps. USA Structures also has its own errors and vintage. The POI reference counts employer establishments at tax addresses and is allocated to tracts by area, so POI validation is approximate. Hazard tables are tract averages, and hidden exposure assumes missing structures share their tract's hazard. The archived 2023 snapshot is a third-party GeoParquet copy of an alpha release. The five regions were chosen by the challenge organisers for their hazards, not to represent the United States.

\section{Conclusion}
A reference-free, conformal completeness model can audit open maps with calibrated uncertainty, detect real losses between releases, and translate gaps into hidden climate-hazard exposure. Applied to Overture Maps, it finds that building gaps in rural and tribal tracts largely closed between 2023 and 2026, that remaining hidden wildfire and drought exposure is concentrated in rural tracts, and that a 2026 release regression removed buildings from high-vulnerability neighbourhoods of San Antonio. It also shows that tract-level flags from completeness models under-detect systematic, group-structured gaps, so equity audits should pair them with reference comparisons aggregated by group.

\paragraph{Data and code availability.} All inputs are public. Code, configuration, the pre-registration and result tables are at \url{https://github.com/GTX-Pradeep/invisible-exposure-open-map-equity}; \texttt{scripts/run\_all.sh} regenerates every number and figure from pinned data versions.

\paragraph{Acknowledgements.} We thank the organisers of the Bias Bounty: Mapping Equity challenge for distributing the strata, and the Overture Maps Foundation, ORNL/FEMA and the U.S. Census Bureau for open data.

\bibliographystyle{plainnat}
\bibliography{refs}

\appendix
\section{Supplementary tables}
\begin{table}[h]\centering\small
\caption{CV scheme comparison (coverage of 90\% ranges, flag rate, log-scale $R^2$).}\label{tab:cv}
\begin{tabular}{lllrrrrr}
\toprule
Layer & Folds & Calibration & Coverage \% & Flags \% & $R^2$ & Region min & Region max \\
\midrule
bldg & random & global & 90.0 & 5.2 & 0.930 & 88.9 & 93.1 \\
bldg & random & mondrian-region & 90.0 & 5.2 & 0.930 & 89.8 & 90.1 \\
bldg & spatial & global & 90.0 & 5.2 & 0.926 & 88.9 & 93.2 \\
bldg & spatial & mondrian-region & 90.0 & 5.2 & 0.926 & 88.9 & 90.2 \\
bldg & region & global & 84.6 & 8.9 & 0.904 & 78.1 & 96.8 \\
bldg & region & mondrian-region & 84.6 & 8.9 & 0.904 & 78.1 & 96.8 \\
poi & random & global & 90.0 & 5.5 & 0.830 & 87.0 & 92.4 \\
poi & random & mondrian-region & 90.1 & 5.5 & 0.830 & 89.9 & 90.2 \\
poi & spatial & global & 90.0 & 5.4 & 0.824 & 85.3 & 93.4 \\
poi & spatial & mondrian-region & 89.9 & 5.4 & 0.824 & 89.7 & 90.1 \\
poi & region & global & 88.1 & 4.9 & 0.806 & 85.4 & 96.3 \\
poi & region & mondrian-region & 88.1 & 4.9 & 0.806 & 85.4 & 96.3 \\
\bottomrule
\end{tabular}
\end{table}
\begin{table}[h]\centering\small
\caption{Exploratory H3 sensitivity: adjusted difference in log-completeness change (region FE, county-clustered 95\% CI).}\label{tab:e3}
\begin{tabular}{llrrrrrr}
\toprule
Variant & Group & Mean (hi) & Mean (lo) & Adj. diff. & CI low & CI high & $p$ \\
\midrule
endpoint 2026-08-19 & svi & 0.107 & 0.063 & 0.041 & 0.000 & 0.082 & 0.049 \\
endpoint 2026-08-19 & cvi & 0.136 & 0.055 & 0.076 & 0.021 & 0.132 & 0.007 \\
endpoint 2026-08-19 & rural & 0.114 & 0.085 & 0.034 & 0.003 & 0.065 & 0.033 \\
endpoint 2026-08-19 & tribal & 0.079 & 0.091 & 0.035 & 0.004 & 0.067 & 0.028 \\
excluding Bexar County & svi & 0.103 & 0.064 & 0.037 & -0.006 & 0.079 & 0.089 \\
excluding Bexar County & cvi & 0.132 & 0.055 & 0.073 & 0.016 & 0.130 & 0.012 \\
excluding Bexar County & rural & 0.119 & 0.083 & 0.040 & 0.008 & 0.072 & 0.013 \\
excluding Bexar County & tribal & 0.079 & 0.090 & 0.036 & 0.002 & 0.070 & 0.039 \\
\bottomrule
\end{tabular}
\end{table}
\begin{table}[h]\centering\small
\caption{Original pre-registered hypotheses (commit \texttt{807831b}), Mann--Whitney U tests as worded. $r_{rb}$: rank-biserial correlation ($a$ vs.\ $b$). Group $a$ is the first-named group.}\label{tab:orig}
{\footnotesize
\begin{tabular}{p{4.3cm}lrrrrrrr}
\toprule
Hypothesis (original wording) & Outcome & $n_a$ & $n_b$ & Mean $a$ & Mean $b$ & $r_{rb}$ & $p$ & $p_{BH}$ \\
\midrule
H1 orig: high-SVI $\times$ high-heat vs low-SVI $\times$ low-heat & he\_per1000 & 1028 & 500 & 4.897 & 2.277 & 0.049 & 0.000 & 0.000 \\
H2 orig: rural/tribal $\times$ high-drought vs non-rural non-tribal (source max share) & b\_src\_max\_share & 464 & 7486 & 0.817 & 0.860 & -0.184 & 0.000 & 0.000 \\
H2 orig: rural/tribal $\times$ high-drought vs non-rural non-tribal (single-source flag) & b\_single\_source & 464 & 7488 & 0.409 & 0.551 & -0.142 & 0.000 & 0.000 \\
H3 orig: shortfall change 2023$\to$2026-10, high- vs low-SVI (log completeness change) & d\_lc & 3362 & 1900 & 0.065 & 0.055 & 0.105 & 0.000 & 0.000 \\
H5 orig: high- vs low-wildfire hazard & he\_per1000 & 2449 & 2449 & 22.179 & 23.104 & -0.010 & 0.103 & 0.103 \\
\bottomrule
\end{tabular}
}
\end{table}
{\footnotesize
\begin{longtable}{lllrrrrr}
\caption{Robustness of building and POI flag contrasts (risk difference, cluster-bootstrap 95\% CI).}\label{tab:robust}\\
\toprule
Variant & Layer & Group & Flags \% & RD & CI low & CI high & aOR \\
\midrule
\endhead
primary & bldg & svi & 5.2 & +0.009 & -0.009 & +0.024 & 1.19 \\
primary & bldg & cvi & 5.2 & +0.018 & -0.005 & +0.036 & 1.25 \\
primary & bldg & rural & 5.2 & -0.002 & -0.022 & +0.018 & 1.06 \\
primary & bldg & tribal & 5.2 & -0.011 & -0.034 & +0.012 & 1.12 \\
primary, excluding HU$<$50 & bldg & svi & 5.1 & +0.010 & -0.009 & +0.025 & 1.21 \\
primary, excluding HU$<$50 & bldg & cvi & 5.1 & +0.018 & -0.005 & +0.035 & 1.26 \\
primary, excluding HU$<$50 & bldg & rural & 5.1 & -0.003 & -0.023 & +0.016 & 1.04 \\
primary, excluding HU$<$50 & bldg & tribal & 5.1 & -0.009 & -0.033 & +0.014 & 1.34 \\
no housing-type covariates & bldg & svi & 6.5 & +0.059 & +0.025 & +0.086 & 2.91 \\
no housing-type covariates & bldg & cvi & 6.5 & +0.038 & -0.001 & +0.074 & 1.65 \\
no housing-type covariates & bldg & rural & 6.5 & -0.021 & -0.045 & +0.007 & 0.73 \\
no housing-type covariates & bldg & tribal & 6.5 & -0.015 & -0.048 & +0.018 & 1.64 \\
negative-binomial GLM & bldg & svi & 5.1 & -0.010 & -0.025 & +0.005 & 0.84 \\
negative-binomial GLM & bldg & cvi & 5.1 & +0.019 & -0.007 & +0.041 & 1.47 \\
negative-binomial GLM & bldg & rural & 5.1 & +0.018 & -0.009 & +0.042 & 1.61 \\
negative-binomial GLM & bldg & tribal & 5.1 & -0.001 & -0.028 & +0.030 & 4.15 \\
80\% range & bldg & svi & 10.3 & +0.009 & -0.014 & +0.030 & 1.10 \\
80\% range & bldg & cvi & 10.3 & +0.027 & -0.004 & +0.054 & 1.21 \\
80\% range & bldg & rural & 10.3 & +0.002 & -0.027 & +0.032 & 1.09 \\
80\% range & bldg & tribal & 10.3 & -0.014 & -0.061 & +0.025 & 1.19 \\
95\% range & bldg & svi & 2.8 & +0.003 & -0.008 & +0.013 & 1.06 \\
95\% range & bldg & cvi & 2.8 & +0.008 & -0.012 & +0.024 & 1.07 \\
95\% range & bldg & rural & 2.8 & -0.005 & -0.019 & +0.009 & 0.94 \\
95\% range & bldg & tribal & 2.8 & -0.012 & -0.025 & +0.001 & 0.85 \\
Mondrian (region) calibration & bldg & svi & 5.2 & +0.006 & -0.012 & +0.020 & 1.12 \\
Mondrian (region) calibration & bldg & cvi & 5.2 & +0.014 & -0.008 & +0.031 & 1.19 \\
Mondrian (region) calibration & bldg & rural & 5.2 & -0.002 & -0.022 & +0.017 & 0.98 \\
Mondrian (region) calibration & bldg & tribal & 5.2 & +0.007 & -0.025 & +0.035 & 1.26 \\
primary & poi & svi & 5.4 & -0.051 & -0.072 & -0.034 & 0.41 \\
primary & poi & cvi & 5.4 & -0.035 & -0.051 & -0.018 & 0.58 \\
primary & poi & rural & 5.4 & -0.021 & -0.035 & -0.009 & 0.56 \\
primary & poi & tribal & 5.4 & -0.005 & -0.026 & +0.012 & 0.84 \\
primary, excluding HU$<$50 & poi & svi & 5.4 & -0.051 & -0.072 & -0.035 & 0.41 \\
primary, excluding HU$<$50 & poi & cvi & 5.4 & -0.034 & -0.051 & -0.018 & 0.57 \\
primary, excluding HU$<$50 & poi & rural & 5.4 & -0.023 & -0.036 & -0.011 & 0.53 \\
primary, excluding HU$<$50 & poi & tribal & 5.4 & -0.005 & -0.025 & +0.013 & 0.83 \\
no housing-type covariates & poi & svi & 5.6 & -0.060 & -0.082 & -0.042 & 0.36 \\
no housing-type covariates & poi & cvi & 5.6 & -0.042 & -0.060 & -0.021 & 0.51 \\
no housing-type covariates & poi & rural & 5.6 & -0.022 & -0.034 & -0.011 & 0.57 \\
no housing-type covariates & poi & tribal & 5.6 & -0.007 & -0.029 & +0.009 & 0.86 \\
negative-binomial GLM & poi & svi & 5.0 & -0.049 & -0.061 & -0.036 & 0.42 \\
negative-binomial GLM & poi & cvi & 5.0 & -0.032 & -0.041 & -0.019 & 0.56 \\
negative-binomial GLM & poi & rural & 5.0 & -0.010 & -0.024 & +0.002 & 0.79 \\
negative-binomial GLM & poi & tribal & 5.0 & -0.008 & -0.026 & +0.007 & 1.67 \\
80\% range & poi & svi & 10.4 & -0.067 & -0.099 & -0.038 & 0.50 \\
80\% range & poi & cvi & 10.4 & -0.053 & -0.084 & -0.019 & 0.67 \\
80\% range & poi & rural & 10.4 & -0.027 & -0.048 & -0.009 & 0.63 \\
80\% range & poi & tribal & 10.4 & +0.003 & -0.029 & +0.026 & 0.89 \\
95\% range & poi & svi & 3.1 & -0.040 & -0.052 & -0.028 & 0.32 \\
95\% range & poi & cvi & 3.1 & -0.025 & -0.037 & -0.010 & 0.47 \\
95\% range & poi & rural & 3.1 & -0.014 & -0.024 & -0.006 & 0.56 \\
95\% range & poi & tribal & 3.1 & -0.012 & -0.028 & -0.002 & 0.79 \\
Mondrian (region) calibration & poi & svi & 5.4 & -0.050 & -0.071 & -0.034 & 0.41 \\
Mondrian (region) calibration & poi & cvi & 5.4 & -0.032 & -0.048 & -0.015 & 0.58 \\
Mondrian (region) calibration & poi & rural & 5.4 & -0.021 & -0.035 & -0.010 & 0.54 \\
Mondrian (region) calibration & poi & tribal & 5.4 & +0.007 & -0.018 & +0.026 & 0.73 \\
reference gap (USA Structures $<$80\%) & bldg & svi & 1.9 & +0.005 & -0.003 & +0.015 & 1.40 \\
reference gap (USA Structures $<$80\%) & bldg & cvi & 1.9 & +0.016 & +0.006 & +0.028 & 2.50 \\
reference gap (USA Structures $<$80\%) & bldg & rural & 1.9 & +0.035 & +0.019 & +0.053 & 3.92 \\
reference gap (USA Structures $<$80\%) & bldg & tribal & 1.9 & +0.027 & +0.002 & +0.069 & 4.19 \\
reference gap (ZBP) & poi & svi & 22.8 & -0.289 & -0.319 & -0.256 & 0.18 \\
reference gap (ZBP) & poi & cvi & 22.8 & -0.180 & -0.207 & -0.146 & 0.35 \\
reference gap (ZBP) & poi & rural & 22.8 & -0.019 & -0.052 & +0.008 & 0.87 \\
reference gap (ZBP) & poi & tribal & 22.8 & +0.019 & -0.020 & +0.059 & 1.20 \\
\bottomrule
\end{longtable}
}
\end{document}
